\chapter{PENDAHULUAN}

\section{Latar Belakang}

% TODO: Tulis latar belakang penelitian Anda di sini.
% Jelaskan konteks, urgensi, dan relevansi topik yang dipilih.
% Rujuk pada penelitian terdahulu dan data sekunder jika ada.

Indonesia memiliki lebih dari 316 ribu sekolah yang terkoneksi internet\cite{indonesiabaikRibuSekolah}.
Dari seluruh sekolah tersebut, seluruhnya diwajibkan untuk mengikuti Tes Kemampuan Akademik (TKA) yang berskala nasional 
dan dilakukan menggunakan sarana komputer atau \emph{Computer-Based Testing} (CBT). Hal ini menimbulkan kebutuhan akan 
kegiatan persiapan ujian TKA bagi tiap sekolah. Kegiatan persiapan ujian TKA harus mensimulasikan kondisi 
ujian TKA yang sesungguhnya, hal ini termasuk akan penggunaan software ujian daring yang mirip.

Platform persiapan ujian daring (Computer-Based Testing/CBT) saat ini menghadapi sejumlah tantangan utama untuk digunakan 
sebagai software persiapan atau simulasi bagi institusi. Dua tantangan utama tersebut yaitu ketersediaan algoritma 
penilaian IRT pada software dan kapasitas software untuk menampung dan mengelola data \emph{concurrent} submision 
dalam jumlah besar atau lebih dari 1000 request per detik. Jumlah data sangat penting bagi institusi karena keakuratan 
parameter IRT sangat bergantung pada jumlah data yang tersedia. Hal ini disebabkan oleh adanya potensi penyimpangan 
statistik dari jumlah data yang tersedia akan mempengaruhi keakuratan parameter IRT sehingga menyebabkan
penurunan akurasi skor.

Di sisi lain, tidak semua institusi pendidikan, sekolah, maupun bimbingan belajar memiliki kapabilitas finansial dan
keahlian rekayasa piranti lunak untuk memelihara infrastruktur mandiri demi satu sistem ujian yang pelaksanaannya
bersifat temporal. Kebutuhan akan solusi \textit{hosted} ini mendasari perlunya sebuah perangkat lunak berbiaya rendah,
mudah diakses, dan dapat diandalkan untuk mengelola ujian daring berskala besar.

Penelitian ini menjadikan platform ujian "Anycademy" sebagai studi kasus teknis \emph{software as a service} multi-tenancy.
Arsitektur \textit{backend} akan dirancang menjadi berbasis microservices dengan modul-modul terpisah meliputi 
autentikasi, manajemen dan pemrosesan data tryout, skoring IRT, manajemen \emph{user}, queue, cache, dan database batching.
Fokus penelitian ini adalah pada modul autentikasi, penilaian berbasis IRT, dan pengelolaan data ujian secara berskala besar
dengan simulasi diatas 1000 peserta secara simultan. 

\section{Rumusan Masalah}

% TODO: Tulis rumusan masalah dalam bentuk pertanyaan penelitian.
% Rumusan masalah harus spesifik, terukur, dan dapat dijawab.

Berdasarkan latar belakang di atas, rumusan masalah dalam penelitian ini adalah:

\begin{enumerate}
	\item Bagaimana merancang arsitektur backend platform ujian daring berbasis SaaS yang memisahkan beban proses bisnis umum secara \textit{modular monolith} dengan pemrosesan matematik IRT ke agen layanan mikro (\textit{microservice}) demi menjaga kestabilan platform?
  \item Bagaimana mengembangkan alur pipa (\textit{pipeline}) pendaftaran nilai berantre dengan bantuan \textit{message broker} guna mengelola penyerahan ujian secara konkuren, guna mencegah \textit{race condition} maupun \textit{downtime}?
	\item Bagaimana performa arsitektur yang dikembangkan saat diukur dan divalidasi terhadap skenario \textit{burst traffic} masif dan simulasi pelanggaran anti-kecurangan berdasarkan metrik \textit{throughput}, latensi, dan utilisasi komputasi?
\end{enumerate}

\section{Batasan Masalah}

% TODO: Tulis batasan masalah untuk memfokuskan penelitian.
% Batasan dapat mencakup ruang lingkup, teknologi, atau aspek yang tidak diteliti.

Agar penelitian ini tetap fokus dan terukur di lingkup Ilmu Komputer, batasan masalah yang ditetapkan adalah:

\begin{enumerate}
	\item \textbf{Fokus Arsitektur Backend:} Penelitian difokuskan pada perancangan arsitektur backend platform Anycademy, yang spesifik mencakup sistem \textit{modular monolith} dengan Go, orkestrasi perpesanan dengan Redis, serta integrasi layanan mikro komputasi berat IRT dengan Python FastAPI.
	\item \textbf{Lingkup Eksklusi Psikometrik:} Analisis terhadap sistem IRT dalam sistem ini murni diukur dari performa teknologi komputasionalnya, seperti Big-O ekseskusi, latensi dan efisiensi memori, tanpa mengukur validitas dan ketepatan perhitungan psikometrik secara keilmuan pendidikan.
	\item \textbf{Implementasi Multi-Tenancy:} Metode \textit{multi-tenancy} dibatasi pada implementasi isolasi data secara logis (\textit{Logical Separation/Row-level atau Schema-level}) pada \textit{database} (PostgreSQL), dan tidak mencakup penerapan pemisahan secara fisik (contohnya pembagian klaster basis data bagi masing-masing institusi).
	\item \textbf{Skala Pengujian Beban:} Skenario pengujian skalabilitas (seperti \textit{Load Testing} dan \textit{Stress Testing}) akan mengutilisasi konfigurasi spesifik untuk skenario pengumpulan ujian (*exam submission*) simulasi serentak.
\end{enumerate}

\section{Tujuan Penelitian}
Tujuan dari penelitian ini adalah:
\begin{enumerate}
	\item Mendesain pola arsitektur backend sistem SaaS yang memisahkan antara beban proses bisnis umum secara \textit{modular monolith} dengan pemrosesan matematik IRT ke agen layanan mikro (*microservice*) demi menjaga kestabilan platform.
	\item Mengembangkan alur pipa (*pipeline*) pendaftaran nilai berantre dengan bantuan \textit{message broker} guna mengelola penyerahan ujian secara konkuren, guna mencegah \textit{race condition} maupun \textit{downtime}.
	\item Mengetahui batas puncak performa server (*hardware limit*) beserta skalabilitas piranti lunaknya lewat analisis statistik kuantitatif terhadap \textit{throughput} (T/s), waktu respons (\textit{response latency}), dan penggunaan CPU/RAM selama simulasi ujian intensif divalidasi.
\end{enumerate}

\section{Manfaat Penelitian}

Hasil dari penelitian ini diharapkan dapat memberikan kontribusi sebagai berikut:

\begin{enumerate}
	\item \textbf{Manfaat Teoritis:} Penelitian ini berkontribusi pada literatur rekayasa perangkat lunak dengan menyajikan
  studi kasus konkret mengenai pola arsitektur hibrida (\textit{polyglot architecture}) untuk domain ed-tech yang 
  menuntut komputasi intensif sekaligus ketahanan terhadap lonjakan trafik konkuren. Secara khusus, penelitian ini 
  memperkaya pemahaman tentang bagaimana pemisahan tanggung jawab antara proses bisnis umum (\textit{modular monolith}) 
  dan komputasi numerik berat (\textit{microservice}) dapat dirumuskan sebagai pola desain yang dapat direplikasi pada 
  sistem serupa.

	\item \textbf{Manfaat Praktis (Industri Perangkat Lunak \& Pengembang EdTech):} Penelitian ini diharapkan dapat menjadi 
  rujukan teknis dan cetak biru (\textit{blueprint}) arsitektur bagi pengembang maupun tim rekayasa yang merancang sistem 
  ujian daring berskala besar, khususnya dalam mengombinasikan keunggulan konkurensi bahasa Go untuk lapisan layanan 
  inti dengan kekuatan ekosistem numerik/ilmiah Python untuk beban komputasi psikometrik. Dengan demikian, hasil penelitian 
  dapat mempercepat proses pengambilan keputusan arsitektural pada proyek-proyek ed-tech serupa, tanpa perlu melalui proses 
  trial-and-error dari nol.

\end{enumerate}


\section{Sistematika Penulisan}

% TODO: Sesuaikan sistematika penulisan dengan struktur penelitian Anda.

Sistematika penulisan proposal tugas akhir ini adalah sebagai berikut:

\begin{itemize}
	\item \textbf{Bab I:} Pendahuluan --- berisi pendefinisian latar belakang problem rekayasa sistem, rumusan arsitektural, ruang lingkup implementasi teknis, luaran fungsionalitas, serta manfaat komputasional program.
	\item \textbf{Bab II:} Tinjauan Pustaka --- merangkum karya-karya terkait \textit{software scaling}, perbandingan arsitektur (\textit{Monolith vs Microservices}), manajemen antrean asinkronus, serta teknologi landasan CBT.
	\item \textbf{Bab III:} Metode Penelitian --- menjabarkan alur metodologi rekayasa perancangan arsitektur, diagaram topologi \textit{SaaS}, pemodelan integrasi modul, skema basis data, dan rancangan pengujian \textit{stress \& load test}.
	\item \textbf{Bab IV:} Jadwal Penelitian --- memuat Gantt Chart linimasa fase perancangan hingga implementasi prototipe.
\end{itemize}
